
\documentclass{beamer}

%\usepackage[table]{xcolor}
\mode<presentation> {
  \usetheme{Boadilla}

%  \usetheme{Pittsburgh}
%\usefonttheme[2]{sans}
\renewcommand{\familydefault}{cmss}
%\usepackage{lmodern}
%\usepackage[T1]{fontenc}
%\usepackage{palatino}
%\usepackage{cmbright}
  \setbeamercovered{transparent}
\useinnertheme{rectangles}
}

\setbeamercolor{normal text}{fg=black}
\setbeamercolor{structure}{fg= blue}
\definecolor{trial}{cmyk}{1,0,0, 0}
\definecolor{trial2}{cmyk}{0.00,0,1, 0}
\definecolor{darkgreen}{rgb}{0,.4, 0.1}
\usepackage{array}
\beamertemplatesolidbackgroundcolor{white}  \setbeamercolor{alerted
text}{fg=red}
\usepackage{tcolorbox}
\setbeamertemplate{caption}[numbered]\newcounter{mylastframe}


\usepackage{tikz}
\usetikzlibrary{arrows}
\usepackage{colortbl}

\renewcommand{\familydefault}{cmss}
%\usepackage[all]{xy}

\usepackage{tikz}
\usepackage{lipsum}

 \newenvironment{changemargin}[3]{%
 \begin{list}{}{%
 \setlength{\topsep}{0pt}%
 \setlength{\leftmargin}{#1}%
 \setlength{\rightmargin}{#2}%
 \setlength{\topmargin}{#3}%
 \setlength{\listparindent}{\parindent}%
 \setlength{\itemindent}{\parindent}%
 \setlength{\parsep}{\parskip}%
 }%
\item[]}{\end{list}}
\usetikzlibrary{arrows}

\usecolortheme{lily}

\newtheorem{com}{Comment}
\newtheorem{lem} {Lemma}
\newtheorem{prop}{Proposition}
\newtheorem{condition}{Condition}
\newtheorem{thm}{Theorem}
\newtheorem{defn}{Definition}
\newtheorem{cor}{Corollary}
\newtheorem{obs}{Observation}
 \numberwithin{equation}{section}
 

\setbeamertemplate{navigation symbols}{}


\title[PLSC 43401]{Mathematical Foundations of Political Methodology}

\author{Robert Gulotty}
\institute[Chicago]{University of Chicago}
\vspace{0.3in}




\begin{document}

\begin{frame}
\maketitle
\end{frame}


\begin{frame}
\frametitle{Day 7: CLT}
\begin{itemize}
\item Moment Generating Functions
\item Limit Theorems
\item Properties of Estimators
\end{itemize}
\end{frame}





\begin{frame}
\frametitle{Approximating functions and second order conditions}


\begin{thm}
\textbf{Taylor's Theorem}
Suppose $f:\mathbb{R} \rightarrow \mathbb{R}$, $f(x)$ is infinitely differentiable function.  Then, the Taylor expansion of $f(x)$ around \alert{$a$} is given by 

\begin{eqnarray}
f(x) & = & f(a) + \frac{f^{'}(a)}{1!} (x- a) + \frac{f^{''}(a)}{2!} (x - a)^2 + \frac{f^{'''}(a)}{3!}(x- a)^3 + \hdots \nonumber \\
f(x) & = & \sum_{n=0}^{\infty } \frac{f^{n} (a) }{n!} (x - a)^n \nonumber
\end{eqnarray}

\end{thm}




\end{frame}


\begin{frame}
\frametitle{Example Function}

Suppose $a$ = 0 and $f(x) = e^{x}$.  Then, 
\begin{eqnarray}
f^{'}(x) & = & e^{x} \nonumber \\
f^{''}(x) & = & e^{x} \nonumber \\
\vdots  & \vdots & \vdots \nonumber \\
f^{n} (x) & = & e^{x} \nonumber 
\end{eqnarray}

This implies
\begin{eqnarray}
e^{x} & = & 1 + x + \frac{x^{2}}{2!} + \frac{x^3}{3!} \hdots + \frac{x^{n}}{n!} + \hdots \nonumber \end{eqnarray}

\end{frame}





\begin{frame}
\frametitle{Moment Generating Functions}


\begin{defn}
Suppose $X$ is a random variable with pdf $f$.  Define, 
\begin{eqnarray}
E[X^{n}] & = & \int_{-\infty}^{\infty} x^{n} f(x) dx   \nonumber 
\end{eqnarray}

We will call $X^{n}$ the \alert{$n^{\text{th}}$} moment of $X$ 

\end{defn}

\begin{itemize}
\item[-] By this definition $var(X) = \text{Second Moment} - \text{First Moment}^{2} $
\item[-] We are assuming that the integral converges
\end{itemize}

\end{frame}


\begin{frame}
\frametitle{Moment Generating Functions}

\begin{prop}
Suppose $X$ is a random variable with pdf $f(x)$.  Call $M(t) = E[e^{tX}]$, 
\begin{eqnarray}
M(t) & = & E[e^{tX}] \nonumber \\
 & = & \int_{-\infty}^{\infty} e^{tx} f(x) dx \nonumber 
 \end{eqnarray}

We will call $M(t)$ the moment generating function, because:
\begin{eqnarray}
\frac{\partial^{n} M (t) }{\partial^{n} t}|_{0} & = & E[X^{n}] \nonumber 
\end{eqnarray}

(Assuming that we can interchange derivative and integral)

\end{prop}



\end{frame}










\begin{frame}
\frametitle{Moment Generating Functions}

\begin{small}
\begin{proof}
Recall the Taylor Expansion of $e^{tX}$ at $0$, \pause 
\begin{eqnarray}
\invisible<1>{e^{tX} & = & 1 + tx + \frac{t^2 x^2}{2!} + \frac{t^3 x^3}{3!} + \hdots \nonumber } \pause 
\end{eqnarray}

\invisible<1-2>{Then, } \pause 
\begin{eqnarray}
\invisible<1-3>{E[e^{tX} ] & = &  1 + tE[X] + \frac{t^2}{2!} E[X^2] + \frac{t^3}{3!} E[X^3] + \hdots \nonumber } \pause 
\end{eqnarray}

\invisible<1-4>{Differentiate once:} \pause 
\begin{eqnarray}
\invisible<1-5>{\frac{\partial M(t)}{\partial t} & = & 0 + E[X]  + \frac{2t}{2!} E[X^2] + \hdots \nonumber \\
M^{'}(0) & = & 0 + E[X] + 0 + 0 \hdots \nonumber } 
\end{eqnarray}





\end{proof}

\end{small}

\end{frame}

\begin{frame}

\begin{small}
\begin{proof}

Differentiate $n$ times \pause 
\begin{eqnarray}
\invisible<1>{\frac{\partial^{n} M(t)}{\partial^{n} t } & = & 0 + 0 + 0 + \hdots + \frac{n \times n-1 \times \hdots 2 \times t^{0} E[X^{n}] }{n!}  + \frac{n!t E[X^{n+1}] }{(n+ 1)! } + \hdots \nonumber \\} \pause 
\invisible<1-2>{& = & \frac{n! E[X^{n}] }{n!}  + \frac{n!t E[X^{n+1}] }{(n+ 1)! } + \hdots  \nonumber } \pause 
\end{eqnarray}

\invisible<1-3>{Evaluated at 0, yields $M^{n}(0)  = E[X^{n}] $} \pause 

\end{proof}

\begin{itemize}
\invisible<1-4>{\item[-] If two random variables, $X$ and $Y$ have the same moment generating functions, then $F_{X}(x) = F_{Y}(y)$ for \alert{almost all} $x$.  } 
\end{itemize}


\end{small}
\end{frame}


\begin{frame}
\frametitle{The Moments of the Normal Distribution}
Suppose $Z \sim N(0,1)$.  \pause 

\begin{eqnarray}
\invisible<1>{E[e^{tX}] & = & \frac{1}{\sqrt{2\pi}} \int_{-\infty}^{\infty} e^{tx} e^{-x^2/2} dx \nonumber } \pause 
\end{eqnarray}

\begin{eqnarray} 
\invisible<1-2>{tx - \frac{1}{2} x^2  & = & tx - \frac{1}{2} x^2 - \frac{1}{2}t^2+\frac{1}{2}t^2 \nonumber } \\\pause 
\invisible<1-3>{ & = & -\frac{1}{2}(x-t)^2 +\frac{1}{2}t^2\nonumber   } \pause 
\end{eqnarray}
\begin{eqnarray} 
\invisible<1-4>{E[e^{tX}] & = & e^{\frac{t^{2}}{2}} \frac{1}{\sqrt{2\pi}} \int_{-\infty}^{\infty} e^{-(x- t)^2/2} dx \nonumber \\} \pause 
 \invisible<1-5>{& = & e^{\frac{t^{2}}{2}} \nonumber } 
\end{eqnarray}


\end{frame}

\begin{frame}
\frametitle{Extracting Moments of the Normal Distribution}
$$M(t)=E[e^{tX}]=e^{\frac{t^{2}}{2}}$$
\pause 
\begin{eqnarray}
\invisible<1>{M^{'}(0) & = &E[X] =  e^{t^2/2} t |_{0} = 0 \nonumber \\} \pause
\invisible<1-2>{M^{''} (0 ) & = & E[X^2] =  e^{t^2/2} (t^2 + 1) |_{0} = 1 \nonumber \\} \pause 
\invisible<1-3>{M^{'''} (0) & = & E[X^3] = e^{t^2/2} t (t^2 + 3) |_{0} = 0 \nonumber \\} \pause 
\invisible<1-4>{M^{''''} (0) & = & E[X^4] = e^{t^2/2} (t^4 + 6t^2 + 3)|_{0} = 3 \nonumber } \\ \pause 
\invisible<1-5>{M^{5} (0) & = & E[X^5] = e^{t^2/2} t (t^4 + 10t^2 + 15)|_{0} = 0 \nonumber } \\ \pause 
\invisible<1-6>{M^{6} (0) & = & E[X^6] = e^{t^2/2} (t^6 + 15t^4 + 45t^2 + 15 )|_{0} = 15 \nonumber } \pause 
\end{eqnarray} 


\end{frame}

\begin{frame}

\begin{prop}
Suppose $X_{i}$ are a sequence of independent random variables.  Define 
\begin{eqnarray}
Y & = & \sum_{i=1}^{N} X_{i} \nonumber 
\end{eqnarray}

Then 
\begin{eqnarray}
M_{Y}(t) & = & \prod_{i=1}^{N} M_{X_{i}}(t) \nonumber 
\end{eqnarray}


\end{prop}

\end{frame}

\begin{frame}

\begin{proof}
\begin{eqnarray}
M_{Y}(t) & = & E[e^{tY}] \nonumber \\ \pause 
\invisible<1>{& = & E[e^{t\sum_{i=1}^{N} X_{i} } ] \nonumber \\ } \pause 
 \invisible<1-2>{& = & E[e^{tX_{1} + tX_{2} + \hdots tX_{N} } ] \nonumber \\} \pause 
  \invisible<1-3>{& = & E[e^{tX_{1} }]E[e^{tX_{2} }] \hdots E[e^{tX_{N} }] \text{ (by independence) } \nonumber \\} \pause 
   \invisible<1-4>{& = & \prod_{i=1}^{N} E[e^{tX_{i}}] \nonumber } 
\end{eqnarray}

\end{proof}

\end{frame}



\begin{frame}
\frametitle{Sequences and Convergence} 

Sequence (refresher):\\
\begin{eqnarray}
\left\{a_{i} \right\}_{i=1}^{\infty} & =  & \left\{a_{1}, a_{2}, a_{3}, \hdots, a_{n}, \hdots, \right\} \nonumber 
\end{eqnarray}


\begin{defn}
We say that the sequence $\left\{a_{i} \right\}_{i=1}^{\infty} $ converges to real number $A$ if for each $\epsilon>0$ there is a positive integer $N$ such that for $n\geq N$, $|a_{n} - A| < \epsilon$
\end{defn}


\end{frame}

\begin{frame}
\frametitle{Sequences and Convergence} 


Sequence of functions: 
\begin{eqnarray}
\left\{ f_{i} \right\}_{i=1}^{\infty}  & = & \left\{f_{1}, f_{2}, f_{3}, \hdots, f_{n}, \hdots, \right\}\nonumber 
\end{eqnarray}

\begin{defn}
Suppose $f_{i}: X \rightarrow \Re$ for all $i$.  Then $\left\{ f_{i} \right\}_{i=1}^{\infty} $ converges \alert{pointwise} to $f$ if, for all $x \in X$ and $\epsilon> 0$, there is an $N$ such that for all $n\geq N$,  
\begin{eqnarray}
|f_{n} (x)  - f(x)|<\epsilon \nonumber 
\end{eqnarray}
\end{defn}

This is as strong of a statement as we're likely to make in statistics \\



\end{frame}





\begin{frame}
\frametitle{Convergence Definitions}

Define $\widehat{\theta}_{n}$ to be estimator for $\theta$ based on $n$ observations.  \pause  \\
\invisible<1>{\alert{Sequence} of estimators: increasing sample size} \pause \\
\begin{eqnarray}
\invisible<1-2>{\left\{\widehat{\theta}_{i}\right\}_{i=1}^{n} &= & \left\{\widehat{\theta}_{1}, \widehat{\theta}_{2},    \widehat{\theta}_{3}, \hdots, \widehat{\theta}_{n} \right\} \nonumber} \pause  
 \end{eqnarray}

\invisible<1-3>{Question: What can we say about $\left\{\widehat{\theta}_{i}\right\}_{i=1}^{n}$ as $n\rightarrow \infty$? } \pause 
\begin{itemize}
\invisible<1-4>{\item[-] What is the probability $\widehat{\theta}_{n}$ differs from $\theta$?  } \pause 
\invisible<1-5>{\item[-] What is the probability $\left\{\widehat{\theta}_{i}\right\}_{i=1}^{n} $ converges to $\theta$? } \pause 
\invisible<1-6>{\item[-] What is sampling distribution of $\widehat{\theta}_{n}$ as $n \rightarrow \infty$ ? } 
\end{itemize}


\end{frame}


\begin{frame}
\frametitle{Convergence in Probability}

\pause 
\invisible<1>{\begin{defn} 
We will say the sequence $\widehat{\theta}_{n}$ converges in probability to $\theta$ (perhaps a non-degenerate RV)  if, 
\begin{eqnarray}
\lim_{n\rightarrow\infty} Prob(|\widehat{\theta}_{n} - \theta | > \epsilon ) = 0 \nonumber 
\end{eqnarray}
For any $\epsilon>0$
\end{defn}} \pause 

\begin{itemize}
\invisible<1-2>{\item[-] $\epsilon$ is a tolerance parameter: how much error around $\theta$?} \pause 
\invisible<1-3>{\item[-] In the limit, convergence in probability implies sampling distribution collapses on a spike at $\theta$} \pause 
\invisible<1-4>{\item[-] $\left\{\widehat{\theta}_{i}\right\}$ need not actually converge to $\theta$, only P($|\theta_{n}$ - $\theta| > \epsilon) = 0 $ } 
\end{itemize}

\end{frame}


\begin{frame}
\frametitle{Example (Cassella and Burger)} 
\pause 
\invisible<1>{Suppose $S \sim $ Uniform(0,1).  Define $X(s) = s$. } \pause \\
\invisible<1-2>{Suppose $X_{n}$ is defined as follows: } \pause 
\begin{eqnarray}
\invisible<1-3>{X_{1}(s) = s + I(s \in [0,1]) & ,&  X_{2} (s)  = s +  I(s \in [0,1/2]) \nonumber \\} \pause
\invisible<1-4>{ X_{3}(s)  = s + I(s \in [1/2, 1])& ,&  X_{4}(s) = s + I(s \in [0,1/3]) \nonumber \\} \pause
\invisible<1-5>{ X_{5} (s)  =  s + I(s \in [1/3,2/3]) & ,&  X_{6}(s) = s + I(s \in [2/3, 1]) \nonumber } \pause
\end{eqnarray}
\invisible<1-6>{Does $X_{n}(s)$ pointwise converge to $X(s)$? \\} \pause
\invisible<1-7>{Does $X_{n}(s) $ converge in probability to $X(s)$?} \pause

\begin{eqnarray}
\invisible<1-8>{P(|X_{n} - X | > \epsilon)  & = &P ( s \in [\mathtt{l}_{n}, \mathtt{u}_{n} ] ) \nonumber } \pause
\end{eqnarray} 

\invisible<1-9>{Length of $[\mathtt{l}_{n}, \mathtt{u}_{n}] \rightarrow 0  \Rightarrow P ( s \in [L_{n}, U_{n} ] )   = 0 $} 


\end{frame}





\begin{frame}
\frametitle{Almost Sure Convergence} 

\pause 
\begin{defn} 
\invisible<1>{We will say the sequence $\widehat{\theta}_{n}$ converges almost surely to $\theta$ if, 
\begin{eqnarray}
Prob(\lim_{n \rightarrow \infty} |\widehat{\theta}_{n} - \theta| > \epsilon ) = 0 \nonumber } \pause 
\end{eqnarray}

\end{defn}


\begin{itemize}
\invisible<1-2>{\item[-] This is stronger than convergence in probablity, it says that sequence converges to $\theta$ (almost everywhere) ) } \pause 
\invisible<1-3>{\item[-] Think about definition of random variable: $\widehat{\theta}_{n}$ is a function from sample space to real line.} \pause 
\invisible<1-4>{\item[-]  Almost sure says that, for all outcomes ($s$)  in  sample space ($S$) $s \in S$, } \pause 
\begin{eqnarray} 
\invisible<1-5>{\widehat{\theta}_{n}(s) & \rightarrow& \theta(s) \nonumber } \pause 
\end{eqnarray}
\invisible<1-6>{\item[] Except for a subset $\mathcal{N} \subset S $ such that $P(\mathcal{N}) = 0 $.  } 
\end{itemize}

\end{frame}

\begin{frame}
\frametitle{Example (Cassella and Berger)} 

Suppose $S \sim $ Uniform(0,1).\\
Suppose $X_{n}$ is define as follows: 
\begin{eqnarray}
X_{1}(s) = s + I(s \in [0,1]) & ,&  X_{2} (s)  = s +  I(s \in [0,1/2]) \nonumber \\
 X_{3}(s)  = s + I(s \in [1/2, 1])& ,&  X_{4}(s) = s + I(s \in [0,1/3]) \nonumber \\
 X_{5} (s)  =  s + I(s \in [1/3,2/3]) & ,&  X_{6}(s) = s + I(s \in [2/3, 1]) \nonumber 
\end{eqnarray}
Does $X_{n}(s) $ converge almost surely to $X(s) = s$?\pause 

\invisible<1>{\alert{No!}: the sequence doesn't converge for each $s$} \pause \\
\invisible<1-2>{For each value of $s$ the sequence varies between $s$ and $s + 1$ infinitely often } 
\end{frame}







\begin{frame}
\frametitle{Convergence in Distribution}
We've talked about $\widehat{\theta}_{n}$'s sampling distribution converging to a normal distribution.\pause   \\
\invisible<1>{This is \alert{convergence in distribution} }\pause 

\begin{defn}
\invisible<1-2>{A series of random variables, $X_1,\ X_2,\ ...$ with cdf $F_{n}(x)$ converges in distribution to random variable $Y$ with cdf $F(x)$ if 
\begin{eqnarray}
\lim_{n\rightarrow \infty} F_{n} (x) - F(x)  = 0 \nonumber 
\end{eqnarray} 
For all $x \in \Re$ where $F(x)$ is continuous.  } \pause 
\end{defn}

\begin{itemize}
\invisible<1-3>{\item[-] Weakest form of convergence almost sure $\rightarrow$ probability $\rightarrow$ distribution } \pause 
\invisible<1-4>{\item[-] Says that cdfs are equal, says nothing about convergence of underlying RV } \pause 
\invisible<1-5>{\item[-] Useful for justifying use of some sampling distributions } 
\end{itemize}

\end{frame}


\begin{frame}
\frametitle{Convergence in Distribution $\not \Rightarrow$ Convergence in Probability}

Define $X \sim N(0,1)$ and each $X_{n} = - X$.  Then:

$X_{n} \sim N(0,1)$ for all $n$ so $X_{n}$ trivially converges to $X$. But, 
\begin{eqnarray}
P(|X_{n} - X| > \epsilon ) & = & P(|X + X| > \epsilon) \nonumber \\
							& = & P( |2X| > \epsilon) \nonumber \\
						    & =  & P(|X| > \epsilon/2) \not \leadsto 0 \nonumber 
\end{eqnarray}						    





\end{frame}


\begin{frame}
\frametitle{Central Limit Theorem}

\begin{prop}
Let $X_{1}$, $X_{2}, \hdots$ be a sequence of independent random variables with mean $\mu$ and variance $\sigma^2$.  Let $X_{i}$ have a cdf $P(X_{i} \leq x) = F(x)$ and moment generating function $M(t) = E[e^{tX_{i}}]$ .  Let $S_{n} = \sum_{i=1}^{n} X_{i}$.  Then 

\begin{eqnarray}
\lim_{n\rightarrow \infty} P\left( \frac{S_{n} - \mu n }{\sigma\sqrt{n}} \leq x \right) &  = &  \frac{1}{\sqrt{2\pi} } \int_{-\infty}^{x} \exp\left( -\frac{z^{2} }{2} \right) dz \nonumber 
\end{eqnarray}


\end{prop}

\pause 
\invisible<1>{Proof plan:}
\begin{itemize}
\invisible<1>{\item[1)] Rely on Fact that convergence of MGFs$\leadsto$ convergence in CDFs}
\invisible<1>{\item[2)] Show that MGFs, in limit, converge on normal MGF}
\end{itemize}


\end{frame}


\begin{frame}

\begin{prop}
Let $F_{n}$ be a sequence of cumulative distribution functions with corresponding moment generating functions $M_{n}$.  $F$ be a cdf with the moment generating functions $M$.  If $\lim_{n\rightarrow \infty} M_{n}(t) \rightarrow M(t)$ for all $t$ in some interval, then $F_{n}(x) \leadsto F(x)$ for all $x$ (when $F$ is continuous).  
\end{prop}

\pause 
\invisible<1>{\begin{prop}
Suppose $\lim_{n\rightarrow \infty} c_{n} \rightarrow c$, then 
\begin{eqnarray}
\lim_{n\rightarrow \infty} \left( 1 + \frac{c_{n}}{n}\right)^{n}  & = & e^{c} \nonumber 
\end{eqnarray}

\end{prop}
}

\pause 

\invisible<1-2>{\begin{prop}
Suppose $M(t)$ is a moment generating function of some random variable $X$.  Then $M(0) = 1$.  
\end{prop}}



\end{frame}




\begin{frame}
\frametitle{Proof of Central Limit Theorem}

\large{Proof.}
Suppose $X_{1}, \hdots, X_{n}$ are i.i.d. variables with $E[X] = 0 $, variance $\sigma^{2}_{x}$, Moment Generating Function (MGF) $M_{x}(s)$. \pause  \\

\invisible<1>{Let $S_{n} = \sum_{i=1}^{n} X_{i}$ and $Z_{n} = \frac{S_{n}}{\sigma_{x} \sqrt{n}}$. } \pause \\
\invisible<1-2>{$M_{S_{n}}  = (M_{x}(s))^{n} $ and $M_{Z_{n}} (s) = \left(M_{x} \left(\frac{s}{\sigma_{x} \sqrt{n}}    \right)    \right)^{n}$} \pause \\
\invisible<1-3>{Using Taylor's Theorem we can write
\begin{eqnarray}
M_{x} (s) & = & M_{x} (0)  + s M^{'}_{x}(0) + \frac{1}{2} s^2 M_{x}^{''}(0) + e_{s} \nonumber 
\end{eqnarray}
$e_{s}/s^2 \rightarrow 0 $ as $s\rightarrow 0$.  } 

\end{frame}


\begin{frame}
\begin{eqnarray}
M_{x} (s) & = & M_{x} (0)  + s M^{'}_{x}(0) + \frac{1}{2} s^2 M_{x}^{''}(0) + e_{s} \nonumber 
\end{eqnarray}
\pause
Filling  in the derivatives of $M_{x}$ we have 
\begin{eqnarray}
M_{x} (s) & = & 1  + 0 + \frac{1}{2} s^2\sigma_{x}^{2}  + \underbrace{e_{s}}_{\text{Goes to zero}} \nonumber 
\end{eqnarray}
\pause
Using the fact that $M_{Z_n}(s)=(M_x(\frac{s}{\sigma_X\sqrt{n}}))^n$
\pause
\begin{eqnarray}
M_{Z_{n}}(s) & = & \left(1 + \frac{1}{2}\left( \frac{s}{\sigma_{x} \sqrt{n}} \right)^{2}\sigma_{x}^{2} \right)^{n} \nonumber \\
& = & \left( 1 + \frac{s^{2}/2}{n} \right)^{n} \nonumber \\
\lim_{n\rightarrow \infty} M_{Z_{n}}(s) & = & e^{\frac{s^2}{2}}\nonumber 
\end{eqnarray}

\end{frame}

\begin{frame}{Estimators}
\begin{itemize}
\item Statistics are functions of random variables.  \pause
\item An \emph{estimator} is statistic that is used to infer the value of an unknown parameter.\pause
\item An estimate is a particular realization of the estimator.\pause
\item Estimators are written with a $\hat{\theta}$.\pause
\item Estimators may provide a single point or an interval.
\end{itemize}
\end{frame}



\begin{frame}{Properties of Estimators}
\begin{itemize}
\item Estimators have distributions, expectations and variances.\pause
\item Does the expected value of the estimator equal the true value of the parameter? \pause \alert{Bias}\pause
\item Does the estimator minimize its variance? \pause \alert{Efficiency}\pause
\item Does the estimator converge in probability to the estimated quantity?\pause \alert{Consistency}\pause
\end{itemize}
\end{frame}

\begin{frame}{Bias}
The bias of an estimator is defined to be
$$\text{Bias}(\hat{\theta})=E(\hat{\theta})-\theta$$.
\begin{itemize}
\item If $E(\hat{\theta})>\theta$, the estimator is biased upward.
\item If $E(\hat{\theta})<\theta$, the estimator is biased downward.
\item If $E(\hat{\theta})=\theta$, the estimator is unbiased.
\end{itemize}

\end{frame}

\begin{frame}{Example of Unbiased Estimator}
\begin{itemize}
\item Sample mean $E(\bar{X})-\mu$
\end{itemize}
\begin{eqnarray}
\invisible<1>{E[\bar{X}_{n}] & = & \frac{1}{n}\sum_{i=1}^{n} E[X_{i}] \nonumber \\ } \pause 
\invisible<1-2>{& = & \frac{1}{n} n \mu  = \mu \nonumber } 
\end{eqnarray}
$$E[\bar{X}_{n}] -\mu=0$$
\end{frame}

\begin{frame}{Example of Bias}
\begin{itemize}
\item Uncorrected sample variance: $E[\hat{\sigma}^2]-\sigma^2$.
\end{itemize}
\scriptsize{
\begin{eqnarray}
\invisible<1>{E[\hat{\sigma}^2] & = & E[\frac{1}{n}\sum_{i=1}^{n}(X_i-\bar{X}_n)^2] \nonumber \\ } \pause
\invisible<1-2>{& = & E[\frac{1}{n}\sum_{i=1}^{n}(X_i^2)-\frac{1}{n}2 \bar{X}_n\sum_{i=1}^{n} X_i +\frac{1}{n}\sum_{i=1}^{n} (\bar{X}_n^2)] \nonumber \\ } \pause 
\invisible<1-3>{& = & E[\frac{1}{n}\sum_{i=1}^{n}(X_i^2)-\frac{1}{n}2 \bar{X}_n n \bar{X}_n +\frac{n}{n}\bar{X}_n^2] \nonumber \\ } \pause 
\invisible<1-4>{& = & E[\frac{1}{n}\sum_{i=1}^{n}(X_i^2)- \bar{X}_n^2] \nonumber \\ } \pause 
\invisible<1-5>{& = & \frac{1}{n}\sum_{i=1}^{n}E[X_i^2]- E[\bar{X}_n^2] \nonumber \\ } \pause 
\invisible<1-6>{& = & \frac{1}{n}\sum_{i=1}^{n}[\sigma^2+\mu^2] - \frac{\sigma^2}{n}-\mu^2 \nonumber \\ } \pause 
\invisible<1-7>{& = &\sigma^2- \frac{\sigma^2}{n} \nonumber } \pause 
\end{eqnarray}
}
\end{frame}

\begin{frame}{Efficiency}
Given two unbiased estimators, we can compare their sampling error.
$$\sigma_{\hat{\theta}}^2= E[(\hat{\theta}-E(\hat{\theta}))^2]$$
\pause
\begin{itemize}
\item Estimator $\hat{\theta}$ is more efficient than $\tilde{\theta}$ if $\sigma_{\hat{\theta}}<\sigma_{\tilde{\theta}}$\pause
\item The estimator with a smaller sampling error is more efficient.\pause
\item If an estimator $\hat{\theta}$ has a smaller standard error than any other unbiased estimator, it is called \emph{efficient}.
\end{itemize}
\end{frame}

\begin{frame}{Example of Efficiency}
\begin{itemize}
\item Sample Mean: $\sigma_{\bar{X}}^2=E((\bar{X}-\mu)^2)$\pause
\begin{eqnarray}
\invisible<1>{\text{var}(\bar{X}_{n}) & = & \frac{1}{n^2} \text{var}(\sum_{i=1}^{n} X_{i}) \nonumber \\ } \pause 
 \invisible<1-2>{&= & \frac{1}{n^2} \sum_{i=1}^{n} \text{var}(X_{i}) \nonumber \\ } \pause 
  \invisible<1-3>{& = & \frac{1}{n^2} n \sigma^2 = \frac{\sigma^2}{n} \nonumber } 
\end{eqnarray}
 $$\sigma_{\bar{X}}^2=\frac{\sigma^2}{n}$$
\end{itemize}
\end{frame}

\begin{frame}{Example of Efficiency}
\begin{itemize}
\item  `First Impression' Estimator: $\mathcal{W}(X_1,\ X_2,\ X_3, ... )=X_1$

\begin{eqnarray}
\invisible<1>{\sigma_{\mathcal{W}}^2&=&E((X_1-\mu)^2)\nonumber \\ } \pause 
\invisible<1-2>{ &=&E[X_1^2]-\mu^2 \nonumber \\ } \pause 
\invisible<1-3>{& = & [\sigma^2+\mu^2]-\mu^2\nonumber \\ }\pause 
\invisible<1-4>{& = & [\sigma^2+\mu^2]-\mu^2\nonumber \\ }\pause
\invisible<1-4>{& = & \sigma^2}\pause
\end{eqnarray}
\item Comparing $\sigma_{\mathcal{W}}^2$ to $\sigma_{\bar{X}}^2$
$$ \sigma^2 >\frac{\sigma^2}{n}$$
\end{itemize}
\end{frame}

\begin{frame}{Example of Efficiency}
\begin{itemize}
\item `Five' Estimator:  $\sigma_{\mathcal{F}}^2=E((5-5)^2)$\pause
$$\sigma_{\mathcal{F}}^2=E(0)=0$$\pause
\item What is the Bias?\pause
$$E(\mathcal{F})-\mu)= 5-\mu$$

\end{itemize}
\end{frame}

\begin{frame}{Mean squared Error}
\begin{itemize}
\item If we have a biased estimator, we can compare sampling error by evaluating \emph{mean squared error}
$$\text{MSE}(\hat{\theta})=E(\hat{\theta}-\theta)^2$$
\item MSE compares bias and variance
$$\text{MSE}(\hat{\theta})=[\text{Bias}(\hat{\theta})]^2+V(\hat{\theta})$$
\end{itemize}
\end{frame}


\begin{frame}{Mean squared Error}
\begin{itemize}
\item If we have a biased estimator, we can compare sampling error by evaluating \emph{mean squared error}
$$\text{MSE}(\hat{\theta})=E(\hat{\theta}-\theta)^2$$
\item MSE compares bias and variance
$$\text{MSE}(\hat{\theta})=[\text{Bias}(\hat{\theta})]^2+V(\hat{\theta})$$
$$\text{Efficiency}((\hat{\theta}_1,\hat{\theta}_2))=\frac{\text{MSE}(\hat{\theta}_2)}{\text{MSE}(\hat{\theta}_1)}$$
\end{itemize}
\end{frame}

\begin{frame}{Design of Studies with a known $\mu$ and $\sigma$}
\begin{itemize}
\item How many draws should we get to make sure that our sample mean is close to the truth?\pause
\item Suppose that we assume that our data is normal, with a known mean $\mu$ and $\sigma^2$\pause
\item What is 
$$P((\bar{X}_n-\mu)/(\sigma/\sqrt{n})<z|\mu,\ \sigma^2, n)$$\pause
\item Fix $\alpha$, and we can calculate $n$.
\end{itemize}
\end{frame}


\begin{frame}{Design of Studies}
\begin{itemize}
\item All we need to do is find the smallest $n$ that produces:
$$P((\bar{X}_n-\mu)/(\sigma/\sqrt{n})<z|\mu,\ \sigma^2, n)\leq \alpha $$
\item And solve for n.
\end{itemize}
\end{frame}

\begin{frame}{Sample size when we know $\mu$ and $\sigma$}
\begin{itemize}
\item Imagine that you are running a survey, and have to pick a
sample size ahead of time such that    
$$P( | \bar X -\mu| < \sigma / 5) =.95$$\pause
\item What sample size would you choose? $\Phi^{-1}(.95)=1.644854$\pause
$$\Phi^{-1}(.975)=1.959$$\pause
$$P( \frac{| \bar X_1  -\mu|}{\frac{\sigma}{\sqrt{n}}} < 1.96) =.95$$\pause
$$P(| \bar X_1  -\mu| < \frac{1.96*\sigma}{\sqrt{n}}) = .95$$\pause
$$\frac{1.96*\sigma}{\sqrt{n}}=\frac{\sigma}{5}$$\pause
$$n=(1.96*5)^2=96.04$$

\end{itemize}
\end{frame}



\begin{frame}{Sample size when we know $\mu$ and $\sigma$}
\begin{itemize}
\item What if you have a random variable with distribution $X\sim N(\mu, 12)$ and you want $|\bar{X}-\mu|<3$ with a probability of .95.
$$P( | \bar X -\mu| <3 )=.95$$\pause
\item What sample size would you choose?
$$P( \frac{| \bar X  -\mu|}{\frac{\sigma}{\sqrt{n}}} < 1.96) =.95$$\pause
$$P(| \bar X  -\mu| < \frac{1.96*12}{\sqrt{n}}) = .95$$\pause
$$\frac{1.96*12}{\sqrt{n}}=3$$\pause
$$n=(1.96*12/3)^2=61.5$$

\end{itemize}
\end{frame}



\begin{frame}{Design of studies}
\begin{itemize}
\item If you are asked, `How large a sample size do I need'
\item Ask:
\item How accurately do you need the answer?
\item What level of confidence do you intend to use?
\item What is your current estimate?
\end{itemize}
\end{frame}

\begin{frame}{Finite Sample Properties of Estimators}
\begin{itemize}
\item If $X_i\sim N(\mu,\ \sigma^2)$ 
$$\bar{X}=\frac{1}{n}\sum_{i=1}^nX_i$$\pause
$$S^2=\frac{1}{n-1}\sum_{i=1}^n(X_i-\bar{X})^2$$\pause
\item $\bar{X}$ and $S^2$ are independent:
$$\frac{\bar{X}-\mu}{\sigma/\sqrt{n}}\sim N(0,1)$$\pause
$$\frac{(n-1)S^2}{\sigma^2}\sim \chi^2_{n-1}$$
\end{itemize}
\end{frame}

\end{document}